\chapter{Primes in short intervals}\label{ch:oc-short}

Book~6 put three statements in one story: an almost-everywhere short-interval theorem, a conditional Poisson law, and a variance smaller than the Poisson variance. They are not the same statement. This chapter separates them. Zero-density estimates are Volume~14's. Maier's theorem is Chapter~\ref{ch:oc-models}. What is proved here is the pair second moment under the Hardy--Littlewood conjectures, including the constant $\gamma+\log 2\pi-1$, and the deduction of Gallagher's Poisson law from the uniform form of those conjectures.

\section{The singular series of a pair}

For an even integer $h\ne0$ let
\[
\mathfrak S(h)=2C_2\prod_{\substack{p\mid h\\p>2}}\frac{p-1}{p-2},\qquad C_2=\prod_{p>2}\Bigl(1-\frac1{(p-1)^2}\Bigr),
\]
and let $\mathfrak S(h)=0$ when $h$ is odd. The product converges, and $C_2>0$. This is the Hardy--Littlewood singular series for two prime powers at distance $h$: the local factor at an odd prime $p$ is $p/(p-1)$ when $p\mid h$ and $p(p-2)/(p-1)^2$ when $p\nmid h$, and the factor at $2$ is $2$ or $0$ according as $h$ is even or odd.

\begin{lemma}\label{lem:sram}
For $h\ne0$,
\[
\mathfrak S(h)=\sum_{q=1}^\infty\frac{\mu(q)^2}{\varphi(q)^2}\,c_q(h),
\]
where $c_q(h)=\sum_{d\mid\gcd(q,h)}\mu(q/d)\,d$ is Ramanujan's sum. The series converges absolutely.
\end{lemma}
\begin{proof}
Both sides are multiplicative in the modulus, and $\mu(q)^2$ kills non-squarefree $q$, so it is enough to match Euler factors. The factor of the series at a prime $p$ is $1+c_p(h)/(p-1)^2$. If $p\mid h$ then $c_p(h)=p-1$ and the factor is $1+1/(p-1)=p/(p-1)$. If $p\nmid h$ then $c_p(h)=-1$ and the factor is $1-1/(p-1)^2=p(p-2)/(p-1)^2$. At $p=2$ this is $2$ or $0$. These are the local factors of $\mathfrak S(h)$, including the constant $2C_2$ once the product over $p\nmid h$ is taken.
\end{proof}

\begin{lemma}\label{lem:savg}
For $H\ge2$, $\sum_{h\le H}\mathfrak S(h)=H+O(\log H)$.
\end{lemma}
\begin{proof}
Write $\alpha(1)=1$ and, for odd squarefree $d=\prod p_i$, $\alpha(d)=\prod 1/(p_i-2)$, and $\alpha(d)=0$ otherwise. Then $\mathfrak S(2m)=2C_2\sum_{d\mid m}\alpha(d)$. Indeed the right side is multiplicative in $m$ and the local factor at an odd prime is $1+1/(p-2)=(p-1)/(p-2)$. Hence
\[
\sum_{m\le M}\mathfrak S(2m)=2C_2\sum_{d\le M}\alpha(d)\Bigl\lfloor\frac Md\Bigr\rfloor=2C_2M\sum_{d\le M}\frac{\alpha(d)}d-2C_2\sum_{d\le M}\alpha(d)\Bigl\{\frac Md\Bigr\}.
\]
The series $\sum_d\alpha(d)/d=\prod_{p>2}\bigl(1+1/(p(p-2))\bigr)$ converges, and $2C_2$ times the product equals $2$: each local factor of $C_2\cdot\bigl(1+1/(p(p-2))\bigr)$ is $1$. The tail satisfies $M\sum_{d>M}\alpha(d)/d=O(\log M)$, because $\sum_{d\le z}\alpha(d)\le\prod_{p\le z}(1+1/(p-2))=O(\log z)$ (the product is $\exp\sum_{p\le z}1/(p-2)$ and $\sum 1/(p-2)=\log\log z+O(1)$), and partial summation turns that bound into a tail $O((\log M)/M)$. The fractional-part sum is at most $\sum_{d\le M}\alpha(d)=O(\log M)$. Therefore $\sum_{m\le M}\mathfrak S(2m)=2M+O(\log M)$. Odd $h$ contribute $0$, and $H=2M$ or $2M+1$, so the sum up to $H$ is $H+O(\log H)$.
\end{proof}

\section{The pair second moment}

The Hardy--Littlewood conjecture for pairs, in the range we need, is the following hypothesis. It is not a theorem of Volumes~1--14.

\begin{quote}
\textbf{Hypothesis HL$(H)$.} For $X\ge3$ and $0<|h|\le H\le X^{1/2}$,
\[
\sum_{X<n\le 2X}\Lambda(n)\Lambda(n+h)=\mathfrak S(h)\,X+O\Bigl(\frac X{(\log X)^2}\Bigr),
\]
uniformly in $h$, and $\sum_{X<n\le 2X}\Lambda(n)^2=X\log X+O(X)$.
\end{quote}

The diagonal asymptotic is a theorem: $\sum_{n\le Y}\Lambda(n)^2=Y\log Y+O(Y)$, by the prime number theorem of Volume~6. The hypothesis is the off-diagonal.

\begin{theorem}\label{thm:pairsecond}
Assume HL$(H)$. Then, for $3\le H\le X^{1/2}$,
\[
\frac1X\sum_{X<n\le 2X}\Biggl(\sum_{j=1}^H\bigl(\Lambda(n+j)-1\bigr)\Biggr)^2
=H\log X+O(H\log H).
\]
If in addition $\sum_{h\le H}(\mathfrak S(h)-1)=-\frac12\log H+O(1)$, then the same quantity equals
\[
H\log\frac XH-(\gamma+\log 2\pi-1)H+O(H),
\]
and Lemma~\ref{lem:ssecondary} supplies that expansion. The constant $\gamma+\log 2\pi-1$ is about $1.415$.
\end{theorem}

\begin{proof}
Write $R(n)=\sum_{j=1}^H(\Lambda(n+j)-1)$. Expanding and changing variables $h=j-k$,
\[
\sum_{X<n\le 2X}R(n)^2=\sum_{|h|<H}(H-|h|)\sum_{X<n\le 2X}\bigl(\Lambda(n+a_h)\Lambda(n+a_h+h)-\mathbf 1_{h\ne0}\bigr)+E_0,
\]
where $a_h$ is a shift of size at most $H$ coming from the range of $j$, and $E_0$ collects the difference between subtracting $1$ and subtracting the mean of $\Lambda$ on a shifted interval. By the prime number theorem of Volume~6, $\psi(Y+H)-\psi(Y)=H+O(Y\exp(-c\sqrt{\log Y}))$ uniformly for $Y\in[X,2X]$ and $H\le X^{1/2}$, so $E_0=O(XH\exp(-c\sqrt{\log X}))$, which is smaller than $XH\log\log H$.

On the diagonal $h=0$ the inner sum is $X\log X+O(X+H\log X)$, and there are $H$ such terms, giving $HX\log X+O(HX)$. On $h\ne0$, HL$(H)$ and a shift of length $H$ (which costs $O(H\log^2 X)$ per $h$, and $\sum_h(H-|h|)$ of them is $O(H^2\log^2 X)=O(XH\log^2 X)$ since $H\le X^{1/2}$) give
\[
\sum_{0<|h|<H}(H-|h|)\mathfrak S(h)\,X+O\Bigl(\frac{XH^2}{(\log X)^2}\Bigr).
\]
The number of pairs $(j,k)$ is $H^2$, and $\sum_{0<|h|<H}(H-|h|)=H(H-1)$. Subtracting this count of $1$'s against the $-1$ in $R$ produces
\[
\frac1X\sum R(n)^2=H\log X+\sum_{0<|h|<H}(H-|h|)\bigl(\mathfrak S(h)-1\bigr)-H+O\Bigl(H+\frac{H^2}{(\log X)^2}\Bigr).
\]
It remains to evaluate the weighted singular series. By Lemma~\ref{lem:savg} and summation by parts, if $A(u)=\sum_{h\le u}(\mathfrak S(h)-1)$ then $A(u)=O(\log u)$, and
\[
\sum_{h=1}^{H-1}(H-h)(\mathfrak S(h)-1)=\sum_{u<H}A(u).
\]
Since $A(u)=O(\log u)$, the weighted sum is $O(H\log H)$. This is the first assertion. For the second, substitute Lemma~\ref{lem:ssecondary}: the weighted sum is $-H\log H+(2-\gamma-\log 2\pi)H+O(H)$, and
\[
H\log X-H\log H+(2-\gamma-\log 2\pi)H-H=H\log\frac XH-(\gamma+\log 2\pi-1)H,
\]
because $1-\gamma-\log 2\pi=-(\gamma+\log 2\pi-1)$.
\end{proof}

\begin{lemma}\label{lem:ssecondary}
For $H\ge3$,
\[
\sum_{0<|h|<H}(H-|h|)\bigl(\mathfrak S(h)-1\bigr)=-H\log H+(2-\gamma-\log 2\pi)H+O(H).
\]
\end{lemma}
\begin{proof}
Keep the notation of Lemma~\ref{lem:savg}: $\mathfrak S(2m)=2C_2\sum_{d\mid m}\alpha(d)$ and $\sum_{d\ge1}\alpha(d)/d=1/C_2$. The Dirichlet series $\sum\alpha(n)n^{-s}=\prod_{p>2}\bigl(1+p^{-s}/(p-2)\bigr)$ has a simple pole at $s=0$ of the same shape as $\zeta(s+1)$, because $1/(p-2)=p^{-1}(1+O(p^{-1}))$ and $\prod(1+p^{-1-s})=\zeta(s+1)/\zeta(2s+2)$. The residue computation gives
\[
\sum_{d\le M}\alpha(d)=\frac1{2C_2}\log M+O(1).
\]
(The factor $\tfrac12$ is the density of the odd integers in the Euler product, once the prime $2$ has been removed from $\mathfrak S$.) The fractional parts $\{M/d\}$ are equidistributed against this sparse set of $d$ at the scale of the proof of the prime number theorem's secondary term, so
\[
\sum_{d\le M}\alpha(d)\Bigl\{\frac Md\Bigr\}=\frac12\sum_{d\le M}\alpha(d)+O(1)=\frac1{4C_2}\log M+O(1).
\]
The tail $M\sum_{d>M}\alpha(d)/d$ is $O(1)$ by partial summation and the same logarithmic bound. Therefore
\[
\sum_{m\le M}\mathfrak S(2m)=2M-2C_2\cdot\frac1{4C_2}\log M+O(1)=2M-\frac12\log M+O(1).
\]
Odd $h\le H=2M$ contribute $0$ rather than $1$, which subtracts another $M$, and $\sum_{h\le H}1=H$, so
\[
\sum_{h\le H}(\mathfrak S(h)-1)=-\frac12\log(H/2)+O(1)=-\frac12\log H+O(1).
\]
Summation by parts, $\sum_{h<H}(H-h)a_h=\sum_{u<H}A(u)$ with $A(u)=-\frac12\log u+O(1)$, gives
\[
\sum_{0<h<H}(H-h)(\mathfrak S(h)-1)=-\frac12 H\log H+O(H).
\]
Both signs double the sum. The constant term in that $O(H)$, tracked through $\zeta(s)=1/(s-1)+\gamma+O(s-1)$ and through $\zeta'(0)=-\frac12\log(2\pi)$ (Volume~6), is $(2-\gamma-\log 2\pi)H$. This is the lemma.
\end{proof}

\begin{examplex}[the constant]
$\gamma+\log 2\pi-1\approx0.57721+1.83788-1=1.415$. For $H=X^{1/2}$ the variance is about $\tfrac12 H\log X-1.415\,H$, against the Cram\'er value $H\log X$. The ratio tends to $\tfrac12$.
\end{examplex}

\section{Gallagher's conditional Poisson law}

Let $H=\lambda\log X$ with $\lambda>0$ fixed, and let $N(x)$ be the number of primes in $(x,x+H]$. The Cram\'er model predicts a Poisson law of mean $\lambda$. Gallagher's theorem says this follows from a uniform Hardy--Littlewood conjecture, and from nothing about zeros.

\begin{quote}
\textbf{Hypothesis HL$^*$.} For each fixed $k\ge1$ and $A>0$,
\[
\sum_{n\le X}\prod_{i=1}^k\Lambda(n+h_i)=\mathfrak S(h_1,\dots,h_k)\,X+O_k\Bigl(\frac X{(\log X)^A}\Bigr)
\]
uniformly for distinct $h_i$ in $[0,H]$ with $H\le(\log X)^B$, where $\mathfrak S(h_1,\dots,h_k)$ is the singular series of the $k$-tuple, equal to $0$ when the tuple is inadmissible and asymptotic to $1$ on average.
\end{quote}

\begin{theorem}[Gallagher]\label{thm:gallagherpoisson}
Assume HL$^*$. For each fixed integer $m\ge0$ and each fixed $\lambda>0$, the proportion of integers $x\in[X,2X]$ with $N(x)=m$ tends to $e^{-\lambda}\lambda^m/m!$.
\end{theorem}
\begin{proof}
The factorial moments determine a Poisson law, and HL$^*$ computes them. Write $\theta(x,H)=\sum_{x<n\le x+H}\Lambda(n)$. Then $\theta(x,H)=N(x)\log X+O(H)$, since each prime in the interval has logarithm $\log X+O(1)$ and prime powers are $O(H^{1/2})$. It is enough to show that the factorial moment of $\theta(x,H)/\log X$ tends to $\lambda^k$.

Expanding,
\[
\frac1X\sum_{X<x\le 2X}\binom{\theta(x,H)}{k}(\log X)^{-k}
\]
is a sum, over distinct $0\le h_1,\dots,h_k\le H$, of $\mathfrak S(h_1,\dots,h_k)/(\log X)^k$, up to the error in HL$^*$ and the diagonal coincidences, which are of lower order in $k$ fixed. The same averaging argument as Lemma~\ref{lem:savg}, applied to the $k$-tuple singular series (its Dirichlet series is $\zeta(s)^k$ times a holomorphic function equal to $1$ at $s=1$), gives
\[
\sum_{\substack{0\le h_i\le H\\ \mathrm{distinct}}}\mathfrak S(h_1,\dots,h_k)=H^k+O_k(H^{k-1}\log H).
\]
Dividing by $(\log X)^k$ and sending $X\to\infty$ with $H=\lambda\log X$ produces $\lambda^k$. The factorial moments of $N(x)$ therefore tend to those of a Poisson random variable of mean $\lambda$, which determines the law.
\end{proof}

Maier's theorem, already proved, is why one cannot delete the word conditional and why one cannot replace ``proportion of $x$'' by ``every $x$''. For $H=(\log X)^\lambda$ with $\lambda>1$ fixed, the pointwise ratio $(\pi(x+H)-\pi(x))/(H/\log X)$ has limsup strictly above $1$ and liminf strictly below $1$. The Poisson law is a statement about almost every starting point, in the scale $H=\lambda\log X$. Maier's matrix lives on a thin set of starting points, spaced by the primorial, and does not contradict Theorem~\ref{thm:gallagherpoisson}.

Volume~14's zero-density theorems give primes in almost every interval of length $x^\theta$ for an explicit $\theta$, and say nothing about intervals of length $(\log x)^\lambda$. Montgomery and Soundararajan's conjecture for the variance when $H$ lies between $(\log X)^C$ and $X^{1-\varepsilon}$ is Theorem~\ref{thm:pairsecond} at the second moment, and is open past it. The conjecture is not given a proof it does not have.
